%==============================================================================
% فصل سوم: مواد و روش کار
% جزئیات باید به‌اندازه‌ای باشد که مطالعه قابل تکرار باشد.
%==============================================================================
\ChapterStart
\chapter{مواد و روش کار}
\label{ch:methods}

\section{نوع پژوهش}
\label{sec:study-design}

این پژوهش یک مطالعه مقطعی- تحلیلی (\lr{Cross-sectional}) است که با هدف تعیین میزان پایبندی به اقدامات پیشگیری ثانویه غیردارویی و عوامل مرتبط با آن در بیماران سکته قلبی، انجام شده است.


\section{جامعه، محیط و نمونه پژوهش}
\label{sec:population}


جامعه مورد مطالعه این پژوهش را کلیه بیماران 18 سال و بالاتر تشکیل می‌دهند که دست‌کم یک ماه از وقوع سکته قلبی (\lr{STEMI} یا \lr{NSTEMI}) آن‌ها گذشته باشد، تحت آنژیوپلاستی اولیه، درمان ترومبولیتیک یا عمل \lr{CABG} قرار گرفته باشند و در نیمه دوم سال 1404 برای ویزیت پیگیری به کلینیک قلب بیمارستان شهید مدنی شهر خرم‌آباد مراجعه کنند.



\section{روش نمونه‌گیری}
\label{sec:sampling}


روش نمونه گیری در این مطالعه، به صورت تمام شماری (\lr{Census})، غیر تصادفی و پیوسته انجام شده است. نمونه ها از بین بیماران مراجعه کننده به کلینیک قلب بیمارستان شهید مدنی شهر خرم آباد در نیمه دوم سال 1404 انتخاب شدند.



\section{حجم نمونه و روش محاسبه آن}
\label{sec:sample-size}

جهت تعیین حجم نمونه در این مطالعه مقطعی، از فرمول کوکران\footnote{\lr{Cochran's Formula}} برای برآورد یک نسبت استفاده شد. با توجه به مطالعات پیشین، میزان پایبندی بیماران پس از انفارکتوس میوکارد به راهکارهای پیشگیری ثانویه غیردارویی (مانند رژیم غذایی و اصلاح سبک زندگی) به‌طور قابل‌توجهی پایین می باشد؛ به همین دلیل، این نسبت ۱۰ درصد $(p = 0.10)$ در نظر گرفته شد\cite{Sibhatu2025CADPrevention,Najafi2016TelephoneFollowUp}. حجم نمونه اولیه با فرض سطح اطمینان 95 درصد و میزان خطای قابل قبول 5 درصد $(d = 0.05)$ برابر با 139 بیمار تعیین گردید.

\[
n = \frac{Z^{2} \cdot p(1-p)}{d^{2}}
  = \frac{(1.96)^{2} \cdot 0.10 \cdot 0.90}{(0.05)^{2}}
  \approx 139
\]

جهت جبران ریزش احتمالی و پرسشنامه‌های ناقص، حدود ۱۰ درصد به این مقدار اضافه شد و حجم نمونه نهایی ۱۵۰ بیمار تعیین گردید. این حجم نمونه علاوه بر برآورد نسبت، توان آماری متوسطی $(0.8)$ جهت تحلیل عوامل مرتبط در آزمون‌های دو متغیره فراهم می‌کند.


\newpage
\section{متغیرها}
\begingroup
\raggedbottom
\scriptsize
\setlength{\tabcolsep}{2pt}
\renewcommand{\arraystretch}{1.28}
\setlength{\LTleft}{\fill}
\setlength{\LTright}{\fill}
\setlength{\LTpre}{0pt}
\setlength{\LTpost}{0pt}
\setlength{\LTcapwidth}{\textwidth}

\begin{longtable}{>{\centering\arraybackslash}p{0.04\textwidth}
                  >{\raggedright\arraybackslash}p{0.11\textwidth}
                  >{\raggedright\arraybackslash}p{0.09\textwidth}
                  >{\raggedright\arraybackslash}p{0.11\textwidth}
                  >{\raggedright\arraybackslash}p{0.27\textwidth}
                  >{\raggedright\arraybackslash}p{0.14\textwidth}
                  >{\raggedright\arraybackslash}p{0.12\textwidth}}
\caption{تعریف و شیوه اندازه‌گیری متغیرهای پژوهش}\label{tab:variables}\\
\toprule
\textbf{ردیف} & \textbf{متغیر} & \textbf{نقش} & \textbf{ماهیت و مقیاس} & \textbf{تعریف عملیاتی} & \textbf{ابزار یا منبع} & \textbf{واحد یا کدگذاری} \\
\midrule
\endfirsthead

\multicolumn{7}{c}{\textbf{ادامه جدول \thetable}}\\
\toprule
\textbf{ردیف} & \textbf{متغیر} & \textbf{نقش} & \textbf{ماهیت و مقیاس} & \textbf{تعریف عملیاتی} & \textbf{ابزار یا منبع} & \textbf{واحد یا کدگذاری} \\
\midrule
\endhead

\midrule
\multicolumn{7}{r}{\footnotesize ادامه در صفحه بعد}\\
\endfoot

\bottomrule
\endlastfoot

۱ & مصرف دخانیات & وابسته & کیفی؛ اسمی & مصرف هرگونه دخانیات، از جمله سیگار و قلیان، در ۳۰ روز اخیر & فرم جمع‌آوری اطلاعات & بله/خیر \\
۲ & فعالیت بدنی منظم & وابسته & کیفی؛ اسمی & انجام حداقل ۱۵۰ دقیقه فعالیت بدنی با شدت متوسط در هفته & فرم جمع‌آوری اطلاعات (پرسشنامه \lr{IPAQ-Short}) & بله/خیر \\
۳ & پایبندی به رژیم غذایی & وابسته & کیفی؛ اسمی & پیروی از الگوی تغذیه‌ای محافظ قلب شامل مصرف کافی میوه، سبزیجات، حبوبات و چربی‌های غیراشباع، مانند روغن زیتون، و محدودیت مصرف سدیم، گوشت قرمز و چربی‌های اشباع & فرم جمع‌آوری اطلاعات (پرسشنامه \lr{MEDAS}) & عدم پایبندی: نمره کمتر از ۹؛ پایبندی: نمره ۹ یا بیشتر \\
۴ & وضعیت روانی (مدیریت استرس) & وابسته & کیفی؛ اسمی & وضعیت روان‌شناختی پایدار همراه با نبود علائم بالینی قابل‌توجه اضطراب و افسردگی که مانع خودمراقبتی و پایبندی به درمان شوند & فرم جمع‌آوری اطلاعات (پرسشنامه \lr{PHQ-4}) & نمره ۳ یا بیشتر: غیرپایبند؛ نمره ۰ تا ۲: پایبند \\
۵ & سن & مستقل & کمی؛ پیوسته & تعداد سال‌های عمر فرد در زمان ارزیابی & فرم جمع‌آوری اطلاعات & سال \\
۶ & جنس & مستقل & کیفی؛ اسمی & هویت جنسی بر اساس فنوتیپ & فرم جمع‌آوری اطلاعات & مرد/زن \\
۷ & تحصیلات & مستقل & کیفی؛ رتبه‌ای & آخرین مدرک تحصیلی فرد & فرم جمع‌آوری اطلاعات & دیپلم یا پایین‌تر/بالاتر از دیپلم \\
۸ & وضعیت تأهل & مستقل & کیفی؛ اسمی & وضعیت فرد از نظر مجرد یا متأهل بودن & فرم جمع‌آوری اطلاعات & مجرد/متأهل \\
۹ & شغل & مستقل & کیفی؛ اسمی & وضعیت اشتغال فرد در زمان ارزیابی & فرم جمع‌آوری اطلاعات & غیرشاغل/شاغل \\
۱۰ & محل سکونت & مستقل & کیفی؛ اسمی & محل سکونت فرد از نظر شهری یا روستایی بودن & فرم جمع‌آوری اطلاعات & روستایی/شهری \\
\pagebreak[4]
۱۱ & بیمه درمانی & مستقل & کیفی؛ اسمی & برخورداری از پوشش بیمه درمانی & فرم جمع‌آوری اطلاعات & ندارد/دارد \\
۱۲ & وضعیت اقتصادی (ادراک‌شده) & مستقل & کیفی؛ رتبه‌ای & ارزیابی فرد از توانایی مالی خود برای تأمین مخارج زندگی & فرم جمع‌آوری اطلاعات & ضعیف؛ متوسط؛ خوب \\
۱۳ & تعداد بیماری‌های همراه & مستقل & کمی؛ گسسته & شمار بیماری‌های همراه ثبت‌شده، از جمله \lr{HTN/DM/CKD/HF} و سایر موارد & فرم جمع‌آوری اطلاعات & عدد \\
۱۴ & سابقه خانوادگی بیماری قلبی & مستقل & کیفی؛ اسمی & سابقه بیماری‌های قلبی‌عروقی در بستگان درجه یک & فرم جمع‌آوری اطلاعات & ندارد/دارد \\
۱۵ & روزهای سپری‌شده از رویداد تا ارزیابی & مستقل & کمی؛ پیوسته & تعداد روزهای بین رویداد یا انجام \lr{PCI} و تاریخ ویزیت & فرم جمع‌آوری اطلاعات & روز \\
۱۶ & شرکت در بازتوانی قلبی & مستقل & کیفی؛ اسمی & سابقه حضور در دوره‌های بازتوانی قلبی & فرم جمع‌آوری اطلاعات & خیر/بله \\
۱۷ & سواد سلامت (\lr{Health literacy}) & مستقل & کیفی؛ اسمی & میزان اطمینان فرد به توانایی خود در تکمیل مستقل فرم‌های پزشکی & فرم جمع‌آوری اطلاعات (\lr{SILS}) & سواد ناکافی: اصلاً، کمی یا تاحدودی؛ سواد کافی: زیاد یا کاملاً \\
۱۸ & درک از بیماری (\lr{Illness Perception}) & مستقل & کمی؛ گسسته & نمره فرد در ابعاد شناختی و احساسی ادراک بیماری & فرم جمع‌آوری اطلاعات (\lr{Brief-IPQ}) & بازه ۰ تا ۸۰ \\
۱۹ & پایبندی دارویی & مستقل & کمی؛ گسسته & نمره خودگزارش‌شده فرد در پرسشنامه چهارگویه‌ای پایبندی دارویی & فرم جمع‌آوری اطلاعات (\lr{MMAS-4}) & بازه ۰ تا ۴؛ نمره صفر نشان‌دهنده پایبندی بالا \\
۲۰ & اقدام درمانی انجام‌شده & مستقل & کیفی؛ اسمی & اقدام درمانی انجام‌شده برای سکته قلبی & فرم جمع‌آوری اطلاعات & \lr{PCI/CABG/}\newline\lr{Medical} \\

\end{longtable}
\endgroup

{\footnotesize
\noindent\textbf{یادداشت:}
اختصارات: \lr{IPAQ-Short}، فرم کوتاه پرسشنامه بین‌المللی فعالیت بدنی؛
\lr{MEDAS}، ابزار سنجش پایبندی به رژیم مدیترانه‌ای؛
\lr{PHQ-4}، پرسشنامه چهارگویه‌ای سلامت روان؛
\lr{SILS}، غربالگر تک‌گویه‌ای سواد سلامت؛
\lr{Brief-IPQ}، پرسشنامه مختصر ادراک بیماری؛
\lr{MMAS-4}، مقیاس چهارگویه‌ای پایبندی دارویی موریسکی؛
\lr{PCI}، مداخله کرونری از راه پوست؛ و
\lr{CABG}، جراحی پیوند بای‌پس عروق کرونر.
\par}



\section{معیارهای ورود و خروج}
\label{sec:eligibility}

\subsection*{معیار ورود}

\begin{itemize}
\item سن $\geq 18$ سال
\item سابقه سکته قلبی (حداقل 1 ماه از تشخیص گذشته باشد)
\item داشتن حداقل یکی از ریسک‌فاکتورهای اصلی قلبی-عروقی (نظیر چاقی، مصرف دخانیات یا سابقه کم‌تحرکی) که نیازمند پایش و مدیریت سبک زندگی باشد.
\item مراجعه به مرکز درمانی مورد مطالعه در بازه زمانی اجرای پژوهش
\item رضایت آگاهانه برای شرکت در مطالعه
\item عدم ابتلا به اختلال شناختی/روانی یا مشکلات ارتباطی که مانع اخذ رضایت معتبر یا پاسخ‌دهی قابل اعتماد می‌شود. (دلیریوم، دمانس، یا ناتوانی شدید در ارتباط کلامی)
\end{itemize}

\subsection*{معیار خروج}

\begin{itemize}
\item وجود بیماری‌های جسمی ناتوان‌کننده (مانند مشکلات ارتوپدی یا سرطان پیشرفته) که عملاً امکان تغییر سبک زندگی (به‌ویژه فعالیت بدنی) را از بیمار سلب کرده باشند.
\item وجود بیماری‌های حاد همزمان (مانند عفونت فعال، نارسایی شدید کلیوی نیازمند دیالیز، یا سرطان‌های پیشرفته) که اولویت درمانی بیمار را از «پیشگیری ثانویه قلبی» تغییر داده باشد.
\item بروز درد حاد قفسه سینه یا علائم نارسایی قلبی جبران‌نشده در زمان مراجعه که نیازمند بستری فوری باشد.
\end{itemize}


\section{ابزار گردآوری اطلاعات}
\label{sec:instruments}


ابزار گردآوری داده‌ها یک فرم جامع و ساختاریافته متشکل از متغیرهای دموگرافیک، بالینی و پرسشنامه‌های استاندارد بومی‌سازی شده بود که روایی و پایایی آن‌ها پیش‌تر در جامعه ایرانی تایید شده بود\cite{Bazzazian2010BriefIPQ,Mehrabi2023Morisky,MahdaviRoshan2018MEDAS,KarimiGhasemabad2021BriefIPQ,VasheghaniFarahani2011IPAQ}. بخش اول فرم به ثبت مشخصات فردی و پزشکی شامل سن، جنسیت، سطح تحصیلات، وضعیت تاهل، وضعیت شغلی، محل سکونت، وضعیت پوشش بیمه، وضعیت اقتصادی ادراک‌شده، بیماری‌های همراه، سابقه خانوادگی بیماری‌های قلبی-عروقی، زمان سپری‌شده از رویداد حاد، نوع اقدام درمانی انجام‌شده و سابقه شرکت در دوره‌های بازتوانی قلبی اختصاص داشت. پایبندی دارویی نیز به عنوان یک متغیر همبسته مهم با استفاده از نسخه ۴ گویه‌ای مقیاس پایبندی دارویی موریسکی ارزیابی شد\cite{Mehrabi2023Morisky}. ارزیابی میزان پایبندی به اقدامات پیشگیری ثانویه غیردارویی بر اساس چهار رکن اصلی توصیه شده در راهنماهای بالینی انجمن قلب آمریکا و کالج کاردیولوژی آمریکا صورت گرفت\cite{Rao2025ACS,Smith2011SecondaryPrevention}. سنجش وضعیت مصرف دخانیات بر پایه خودگزارش‌دهی بیماران انجام گرفت و قطع کامل هرگونه مصرف در ۳۰ روز اخیر به عنوان معیار پایبندی لحاظ شد\cite{Smith2011SecondaryPrevention,Hayrumyan2023SmokingCessation,Yudi2017SmokingCessation}. ارزیابی فعالیت بدنی با استفاده از نسخه کوتاه پرسشنامه بین‌المللی فعالیت بدنی\footnote{\lr{The International Physical Activity Questionnaire - Short Form}} (\lr{IPAQ-SF}) انجام شد که بر اساس آن، کسب حداقل ۱۵۰ دقیقه فعالیت بدنی با شدت متوسط در هفته یا معادل آن، مرز ورود به گروه پایبند بود\cite{VasheghaniFarahani2011IPAQ,Smith2011SecondaryPrevention}. الگوی تغذیه‌ای بیماران به وسیله پرسشنامه ۱۴ گویه‌ای غربالگری پایبندی به رژیم مدیترانه‌ای\footnote{\lr{Mediterranean Diet Adherence Screener}} (\lr{MEDAS}) سنجیده شد که در آن کسب نمره ۹ یا بالاتر نشان‌دهنده پایبندی به رژیم غذایی محافظ قلب بود\cite{MahdaviRoshan2018MEDAS,HutchinsWiese2022Mediterranean}. وضعیت مدیریت استرس و سلامت روان نیز توسط پرسشنامه استاندارد چهار سوالی سلامت بیمار\footnote{\lr{Patient Health Questionnaire-4}} (\lr{PHQ-4}) غربالگری شد که نمرات کمتر از ۳ به عنوان وضعیت روانی پایدار و کنترل‌شده در نظر گرفته شدند\cite{Dadfar2025PHQ4,Ahmadi2020PHQ4,Ghaheri2020PHQ4}.



\section{روش گردآوری اطلاعات و اجرای مداخله}
\label{sec:procedure}


پژوهش حاضر یک مطالعه مقطعی-تحلیلی می باشد که با هدف تعیین میزان پایبندی به اقدامات پیشگیری ثانویه غیردارویی و عوامل مرتبط با آن در بیماران مبتلا به سکته قلبی طراحی و اجرا گردید.

جامعه آماری این تحقیق را کلیه بیماران ۱۸ سال و بالاتر تشکیل می‌دادند که دست‌کم یک ماه از وقوع سکته قلبی آن‌ها (شامل \lr{STEMI} یا \lr{NSTEMI}) گذشته بود، تحت مداخلات درمانی نظیر آنژیوپلاستی اولیه (\lr{PPCI})، درمان ترومبولیتیک یا عمل جراحی بای‌پاس عروق کرونر (\lr{CABG}) قرار گرفته بودند و در نیمسال دوم سال ۱۴۰۴ جهت پیگیری، ویزیت دوره ای و درمان به درمانگاه تخصصی قلب شهید مدنی شهر خرم‌آباد مراجعه کرده یا در فرآیند پیگیری این مرکز قرار داشتند.

به منظور برآورد دقيق حجم نمونه، با استناد به فرمول استاندارد محاسبه حجم نمونه برای مطالعات مقطعی برآورد نسبت (فرمول کوکران) و بر اساس مطالعات پیشین\cite{Sibhatu2025CADPrevention,Najafi2016TelephoneFollowUp}، با در نظر گرفتن سطح اطمینان ۹۵ درصد و حاشیه خطای قابل قبول 5 درصد، حجم نمونه مورد نیاز برابر با 150 بیمار تعیین شد. با توجه به اهداف پژوهش و محدودیت حجم جامعه در بازه زمانی تعیین‌شده، روش نمونه‌گیری به صورت تمام‌شماری (\lr{Census}) انتخاب شد. بدین ترتیب، کلیه بیماران واجد شرایط مراجعه‌کننده در این بازه ۶ ماهه وارد فرآیند ارزیابی شدند تا بالاترین توان آماری ممکن برای سنجش حجم اثر متغیرها تأمین گردد.

روند اجرای پژوهش بدین صورت بود که جهت افزایش جامعیت مطالعه، دسترسی به حداکثر تعداد بیماران واجد شرایط و همچنین به حداقل رساندن سوگیری انتخاب (\lr{Selection Bias})، فرآیند جمع‌آوری داده‌ها به دو شیوه مصاحبه حضوری در درمانگاه و مصاحبه تلفنی سازمان‌دهی شد. در گام نخست، پرونده‌های پزشکی و سیستم نوبت‌دهی درمانگاه قلب شهید مدنی مورد ارزیابی دقيق قرار گرفت و بیمارانی که تشخیص قطعی سکته قلبی را دریافت کرده بودند، به عنوان کاندید ورود به مطالعه در نظر گرفته شدند. در شیوه حضوری، ارزیابی بیمار در زمان مراجعه به کلینیک انجام شد. در شیوه تلفنی، با بیمارانی که به دلایل مختلف در بازه زمانی مورد نظر به صورت حضوری به کلینیک مراجعه نکرده بودند، تماس تلفنی برقرار گردید. در هر دو روش، ابتدا اهداف پژوهش به طور کامل برای بیماران تشریح شدند و پس از بررسی معیارهای واجد شرایط بودن و کسب رضایت آگاهانه، پرسش‌های مطالعه مطرح گردیدند.

معیار های اصلی ورود بیماران به مطالعه شامل سن برابر یا بیش از ۱۸ سال، گذشت حداقل یک ماه از ابتلا به سکته قلبی و داشتن پرونده پزشکی معتبر در مرکز فوق بود. در سوی مقابل، وجود بیماری‌های جسمی ناتوان‌کننده نظیر مشکلات شدید ارتوپدی یا بدخیمی‌های پیشرفته که عملاً امکان تغییر سبک زندگی و انجام فعالیت بدنی را از بیمار سلب می‌کرد، ابتلا به بیماری‌های حاد همزمان مانند عفونت فعال یا نارسایی شدید کلیوی نیازمند دیالیز که اولویت درمانی را تغییر می‌داد و بروز درد حاد قفسه سینه یا علائم نارسایی قلبی جبران‌نشده در زمان مراجعه که نیازمند بستری فوری بود، به عنوان معیارهای خروج از پیش تعیین‌شده در نظر گرفته شدند.


\section{روش‌های آماری و تجزیه‌وتحلیل داده‌ها}
\label{sec:statistics}


پس از گردآوری داده‌ها، اطلاعات استخراج شده وارد نرم‌افزار آماری \lr{SPSS} (نسخه 27) شده و مورد تجزيه و تحلیل قرار گرفتند. در بخش آمار توصیفی، از شاخص‌های میانگین و انحراف معیار برای متغیرهای کمی، و از شاخص‌های فراوانی مطلق و درصدی برای متغیرهای کیفی استفاده شد. پیش‌فرض نرمال بودن توزیع داده‌های کمی ابتدا توسط آزمون کلموگروف-اسمیرنوف مورد ارزیابی قرار گرفت. به منظور بررسی ارتباط و مقایسه میزان پایبندی به اقدامات غیردارویی در میان گروه‌های مختلف دموگرافیک و بالینی، بر حسب نوع و توزیع متغیرها، از آزمون‌های مجذور کای (\lr{Chi-Square})، آزمون دقیق فیشر، آزمون تی مستقل (\lr{Independent t-test}) یا معادل ناپارامتری آن (من-ویتنی) استفاده شد. در نهایت، به منظور حذف اثر متغیرهای مخدوش‌کننده و شناسایی عوامل پیش‌بینی‌کننده مستقل و همزمان پایبندی غیردارویی، مدل رگرسیون لجستیک باینری (\lr{Binary Logistic Regression}) اجرا گردید. در تمامی تحلیل‌ها، سطح معنی‌داری آماری کمتر از $0.05$ $(P < 0.05)$ ملاک تصمیم‌گیری قرار گرفت.

\section{محدودیت های پژوهش}


این مطالعه به‌صورت تک‌مرکزی در درمانگاه قلب بیمارستان شهید مدنی انجام شده است. بنابراین حجم نمونه و نتایج به بیماران این مرکز محدود است و تعمیم آن به سایر بیمارستان‌ها یا جمعیت‌های روستایی/شهری دیگر می بایست با احتیاط صورت ‌گیرد. همچنین طراحی این مطالعه از نوع مقطعی می باشد و داده‌ها در یک زمان جمع‌آوری شدند لذا امکان بررسی روند زمانی پایبندی و استنتاج روابط علت و معلولی را فراهم نمی‌سازد، بلکه تنها ارتباط بین متغیرها را نشان می‌دهد. به‌دلیل محدودیت زمانی و ظرفیت مرکز، حجم نمونه نسبتا محدود می باشد و این موضوع ممکن است باعث کاهش توان مطالعه برای شناسایی \lr{small effect sizes} شود؛ هرچند توان مطالعه برای بررسی ارتباطات در حد \lr{medium effect size} به‌طور مناسب تأمین خواهد شد.

نمونه‌گیری به‌صورت متوالی از مراجعه‌کنندگان به درمانگاه انجام شده است و احتمال سوگیری انتخاب (\lr{selection/sampling bias}) وجود دارد. به‌ویژه بیمارانی که برای پیگیری مراجعه نمی‌کنند یا دسترسی دشوارتری دارند ممکن است کمتر وارد نمونه شده باشند و نتایج ممکن است پایبندی واقعی در جمعیت عمومی بیماران سکته قلبی را بیش‌ازحد برآورد کند. برای کاهش این سوگیری، فرایند غربالگری در تمام روزهای کلینیک و طی ساعت‌های متعدد انجام شده است. همچنین استفاده از پرسشنامه جهت جمع آوری اطلاعات، احتمال سوگیری یادآوری (\lr{recall bias}) و تمایل به پاسخ‌های مطلوب اجتماعی (\lr{social desirability bias}) را افزایش می دهد. بی نام بودن پرسشنامه، تأکید بر محرمانگی و استفاده از ابزار استاندارد در کاهش این سوگیری ها تاثیرگذار می باشند.


\section{ملاحظات اخلاقی}
\label{sec:ethics}


تمام مراحل اجرای این پژوهش با رعایت کامل موازین اخلاقی مصوب دانشگاه علوم پزشکی لرستان، از جمله حفظ محرمانگی کامل اطلاعات هویت بیماران، عدم استفاده از نام مراجعان در فایلهای تحلیلی، آزادی کامل برای خروج از مطالعه در هر مرحله، حصول اطمینان از عدم تاثیر شرکت در طرح بر روند مراقبت‌های درمانی روتین بیماران و رعایت موارد زیر انجام گرفت:

\begin{itemize}
\item كسب مجوز انجام پژوهش از کمیته اخلاق دانشگاه علوم پزشکی لرستان
\item کسب اجازه از ریاست و معاونت پژوهشی دانشگاه علوم پزشکی لرستان
\item معرفی خود به واحدهای پژوهش و توضیح اهداف پژوهش
\item ارائه نتایج به مراکز موردپژوهش در صورت درخواست
\item محرمانه ماندن اطلاعات و استفاده از اطلاعات فقط جهت تجزیه‌وتحلیل آماری
\end{itemize}
